\section{Mechanistic interpretation and key questions}
\label{app:clarifications}

\noindent\textbf{Q1. What do the response law and decompositions in \cref{sec:loop,sec:local,sec:broadcast} explain beyond their familiar mathematical ingredients?}
\par\noindent
They provide a connected mechanistic account of how dense local supervision can produce either weak transfer or deteriorating generation. OPD updates the policy that generates its next training population. Understanding its dynamics therefore requires following supervision beyond the current loss, through its effects on future behavior and training histories (student's rollouts and the induced supervision signals / request field).

The theory separates three links in this process. First, \cref{sec:local} explains how a teacher's requests become parameter motion: accessibility determines their expression through the student's tangent map, while coherence determines how the resulting gradients combine. Teacher benchmark performance, magnitude of (KL-based) disagreement, or gradient norm alone cannot characterize this conversion. Second, \cref{sec:broadcast} distinguishes changing a prediction from changing the ancestral decisions that make its history likely. Shared-parameter broadcast can increase exposure to a continuation without making the student better at executing it; the resulting histories then change the supervision received from the same fixed teacher. Third, \cref{eq:broadcast-continuations} connects the delivered response to completed outcomes through what the \emph{student's remaining policy} can achieve, separating teacher matching from useful reasoning.

Together, these links explain the structure of the empirical puzzles. The stronger-teacher study exposes the importance of gradient coherence and tangent accessibility; the degenerate run exposes that, under substantial policy motion, the degree of freedom granted to the learner to generate its own rollouts (training histories) can make it generate the new rollouts (reasoning states) which it has difficulty managing and grows uncertain on, inducing poorer-quality supervision signals. Loss reduction dominated by local fitting on these poorly supervised rollouts cannot improve the learner's terminal scores. The mechanisms revealed in the degenerate run motivates preparing the student for the continuations that OPD will elicit, even when preparation does not substantially improve current accuracy by itself. The theoretical contribution is the integrated explanation that makes these mechanisms distinguishable, connects them to outcomes, and motivates an intervention that boosts the model's learning stability and generalization.

\Needspace{5\baselineskip}
\noindent\textbf{Q2. Why is every training run reported with a single seed?}

\textbf{Cost}. The runs behind this paper were trained and evaluated on GH200 nodes, and the single-seed set already consumed roughly 4{,}000 GPU-hours, which is at the limit of our computing resources. OPD is known to be computationally demanding, we therefore designed the experiments so that the main comparisons do not rest on seed variance. The mechanism measurements of Sections 3.2 and 3.3 are exact identities evaluated within a run, on fixed receiver banks and fixed checkpoints, and their uncertainty is reported at the level of problems or sites (paired bootstrap intervals). The outcome comparisons of Section 3.4 are paired within a student-teacher pair on the same problems and replicated across three pairs and two student scales. What single seeds leave unquantified is run-to-run variation of the training trajectories themselves.

\Needspace{5\baselineskip}
\noindent\textbf{Q3. Why does \cref{sec:broadcast} need both prefix-level routing/drift (\cref{eq:rd-law}) and trajectory-level continuation advantages (\cref{eq:broadcast-continuations}) to explain degeneration?}
\par\noindent
For loss or entropy, $J_Y=\mathbb{E}_{\nu_t}Y_t(x)$ describes predictions on the currently observed histories. Cross-scoring separates changes in those histories from changes in predictions on the \emph{same} histories. In \cref{fig:rd-contrast}, the initial student is already highly uncertain on the deep histories that later become prevalent, while subsequent fitting sharpens them. On identical batch-41 histories, discrepancy falls $0.79\to0.06$ and entropy $4.99\to3.99$. Together with the stopping/reach probes, these measurements distinguish routing followed by fitting from uniform entropy-increasing drift (\cref{app:rd-failure,app:rd-entropy-evidence}).

But fitting a prefix and improving the completion from it are different effects. For terminal correctness, $J_Y=\mathbb{E}_{\Pi_t}Y$, and the conditional value $V_t^Y(x)$ depends on \emph{every downstream decision}. At an intermediate cut of the autoregressive trace, its derivative separates changes in reaching the cut from changes in the remaining policy. Unrolling the latter gives the continuation-advantage weighting in \cref{eq:broadcast-continuations}; replacing it with a derivative of the logits at $x$ alone would omit downstream effects (\cref{app:outcome-proof}). Thus, a teacher-preferred next token need not help if the student cannot complete the reasoning it initiates. The prefix analysis establishes the observed routing-then-fitting pattern; the outcome analysis explains why local fitting need not restore correctness. Entropy enters the argument only as a prefix-level readout.

\Needspace{5\baselineskip}
\noindent\textbf{Q4. Could the responsiveness changes in \cref{sec:local} (Figure~\ref{fig:collective-dynamics}) simply reflect changing teacher--student disagreement or generic kernel movement?}
\par\noindent
A live accessibility measure $A_t[B_t]$ can change because either the requests or the kernel change. We therefore fix the initial numerical requests $B_0$, exact histories, and output coordinates across checkpoints. Since only the kernel varies in $A_t[B_0]$, its change measures responsiveness to the \emph{same supervision} (\cref{app:collective-frozen}).

To test whether eigenspace movement reflects teacher-specific adaptation, we evaluate both teachers' frozen fields in the \emph{same checkpoint kernel} and compare their projection gains, separating preferential reorientation from shared movement that benefits both fields. Healthy 1.5B exhibits preferential reorientation; healthy 7B mainly amplifies existing directions; stalled runs lose sensitivity to the fixed requests (\cref{app:collective-specificity}). These findings complement the initial accessibility/coherence decomposition, which exposes conversion differences concealed by comparable gradient norms across the runs.

\Needspace{5\baselineskip}
\noindent\textbf{Q5. What does preparation in \cref{sec:preparation,fig:preparation-aime} contribute beyond offline warm-up, and how does \cref{sec:broadcast} motivate it?}
\par\noindent
The contribution is an explanation of why the value of a warm-up stage, already used in prior work (\cref{app:related-work}), should be assessed through \emph{subsequent learning}. The failure mechanism identifies a mismatch between making extended reasoning likely and being able to execute it. Offline distillation supplies teacher-generated reasoning histories without requiring the initial student to generate them first. It can therefore prepare continuation skills and reasoning-state management that become useful as OPD changes visitation: \emph{continuation readiness}, a benefit beyond higher starting accuracy -- our pass@k results clearly demonstrate this.

